\begin{document}
\maketitle

\begin{abstract}
Let $(Q_j)$ be the denominators of the convergents of $e$ that immediately
precede the partial quotients $2j$, so that $Q_{j+1}=(4j+2)Q_j+Q_{j-1}$. We
study the weighted approximation function
$\Lambda(x)=\limsup_{j\to\infty}j\|Q_jx\|$ and its level sets. The spectrum
is $\Lambda(\R)=\frac14\Z_{\ge0}\cup\{\infty\}$, and $\Lambda$ vanishes exactly
on the group $L=\Z+\Z\frac{e+1}{2}$. Every finite sublevel set
$\{\Lambda\le c\}$ has Hausdorff dimension $0$, whereas for the weights
$j^\alpha$ with $0\le\alpha<1$ the sublevel and prescribed-limit sets have
dimension $1-\alpha$. At each level $k/4$ we determine the exact Hausdorff
gauge. For $h_k(1/Q_n)=(2k+1)^{-n}$ the level set has infinite but
$\sigma$-finite $\mathcal H^{h_k}$-measure in every open arc of $\R/\Z$ and the
strict sublevel is $\mathcal H^{h_k}$-null; every gauge is classified by its
lower limit along the scales $1/Q_n$. Each positive level moreover splits
three ways. The set where $j\|Q_jx\|\to k/4$ has critical gauge base $2$ for
every $k$, and each set where $j\,(Q_jx-\mathrm{round}(Q_jx))\to\pm k/4$ is a
single coset of $L$. The proofs rest on an integer digit expansion attached to
the three-term recurrence.
\end{abstract}

\medskip
\noindent\textit{2020 Mathematics Subject Classification.} 11J83, 11J70, 28A78.

\noindent\textit{Keywords.} Metric Diophantine approximation, continued
fraction of $e$, approximation spectrum, Hausdorff dimension, Hausdorff gauge.
